\begin{textsample}{POS Dim 2 – ro – Score 24.00 – ro/ro\_em\_5.txt}  \label{ex:f2_pos_189}
1.3.1.
Contexto Estadual O Estado de Rondônia vivencia um círculo virtuoso de crescimento e desenvolvimento, cuja sustentabilidade representa um grande desafio e impõe diversas demandas.
Dentre as quais, evidencia -se claramente a necessidade de intensificar investimentos na expansão do Ensino Médio com \textbf{qualidade}.
Estudos realizados por Ramos \& Gomes ( 2015 ) sobre estimativas de demandas educacionais no Ensino Médio no Plano Estadual de Educação ( PEE ) apontam para cenários que “ necessitam não só de aportes de recursos financeiros e humanos para ampliar a \textbf{oferta} de vagas àqueles que venham requerer, mas também, o \textbf{fortalecimento} institucional na execução das \textbf{políticas} públicas traçadas para o Ensino Médio ” ( RONDÔNIA, 2014, p. 40 ).
Rondônia possui diversos \textbf{distritos}, reservas extrativistas, assentamentos de \textbf{reforma} agrária e comunidades indígenas que vivem às margens dos rios e espalhadas por seu território ( BRASIL, 2016 ).
Devido a essas características, que tornam as localidades do interior em difícil acesso, o ensino médio não pode ser \textbf{ofertado} para considerável parcela dos jovens das áreas rurais.
Assim, em decorrência de tais aspectos, muitos migram do campo para as cidades na esperança de continuar os estudos e, consequentemente, concluir a Educação Básica.
Contudo, devido às dificuldades para se manterem nos centros urbanos, uma grande quantidade abandona a escola.
Mediante tal panorama, esses aspectos impulsionam o Estado para que busque distintas possibilidades de atendimento educacional, conforme a diversidade dos cenários.
Assim, a implementação de uma \textbf{política} pública que \textbf{atenda} aos estudantes do ensino médio, em cada uma das suas especificidades, se apresenta de forma relevante.
A universalização do Ensino Médio é assunto de destaque na pauta das discussões educacionais contemporâneas, e representa um desafio, principalmente considerando o número de jovens com distorção idade/série e fora dessa etapa de ensino, em decorrência da \textbf{reprovação} e do abandono da escola.
Diante desse cenário, vários fatores corroboram para o abandono dos estudos, como a exigência de trabalho precoce, ambiente educacional desinteressante, dificuldades no processo ensino-aprendizagem, falta de estímulo dos familiares, difícil acesso ao ambiente escolar, distanciamento entre os saberes escolares e as situações vivenciadas na experiência cotidiana dos jovens, entre outros.
Em face desse contexto, há a necessidade de \textbf{oferta} e ampliação do atendimento escolar para população de 15 a 17 anos, e de elevar a taxa líquida de matrículas do Ensino Médio nessa faixa etária até 2024.
Tornando o atendimento à meta 3 do Plano Nacional de Educação ( PNE ) – relacionada à universalização do atendimento escolar, um grande desafio e prioridade na agenda governamental do Estado de Rondônia ( RONDÔNIA, 2016 ).
A Secretaria de Estado da Educação, por meio do Governo do Estado de Rondônia, almeja atingir as metas do PNE ( 2014-2024 ) que assinalam para a necessidade de, até o final da vigência do plano, aumentar “ [... ] a taxa líquida de matrículas no ensino médio para 85 \% ” ( BRASIL, 2014, p.62 ).
Uma forma de avaliação da \textbf{rede} escolar é o cálculo do Índice de Desenvolvimento da Educação Básica, o IDEB, que considera o fluxo escolar e o desempenho dos estudantes em avaliações padronizadas.
O índice, que varia de 0 a 10, combina dados de aprovação do \textbf{Censo} Escolar com as pontuações obtidas no Sistema de Avaliação da Educação Básica ( SAEB ).
A metodologia consiste, portanto, em um entrecruzamento dos índices de fluxo e de aprendizagem.
O IDEB de Rondônia saiu de 2,96 em 2005 para 3,76 em 2017, ultrapassando assim a \textbf{média} nacional de 3,47.
No entanto, vale pontuar que essa \textbf{trajetória} apresentou períodos de queda e recuperação da \textbf{média} \textbf{estadual}, a qual, em linhas gerais, esteve bastante próxima aos resultados apresentados pelo país.
Em 2009 e 2017, o estado alcançou seus melhores resultados, 3,69 e 3,76, respectivamente, em ambos os casos ultrapassando inclusive o desempenho do país, em uma \textbf{média} de 0,31 pontos.
Destaca -se, por fim, que, em 2019, as escolas públicas \textbf{estaduais} atingiram a \textbf{média} 4,0, mas ainda não atingiu a meta que é de 4,5, contudo no último intervalo de medição, entre 2017 e 2019, o estado obteve uma melhora de 0,2 pontos em seu resultado. ( RONDÔNIA, 2014-2024 ).
Mediante os indicadores relacionados às avaliações externas e internas, foi iniciado o processo de reformulação dessa etapa de ensino para garantir o direito de aprender aos jovens de 15 a 17 anos, propondo intervenções no sistema de ensino \textbf{estadual} para torná -lo mais inclusivo e \textbf{equânime}.
Almejando, assim, possibilitar aos jovens o acesso, permanência e sucesso escolar, com aprendizagem dos saberes culturais e científicos para o exercício da cidadania e inserção no mundo do trabalho.
Dessa forma, busca -se atingir patamares mais elevados na universalização do atendimento e promoção do sucesso escolar do estudante do Ensino Médio.
Com o intuito de alcançar as \textbf{médias}, haverá o intermédio da implementação da BNCC, a ampliação da \textbf{carga} \textbf{horária} e a flexibilização do \textbf{currículo} por meio dos \textbf{Itinerários} \textbf{Formativos} por área do conhecimento e/ou \textbf{Itinerários} de Formação \textbf{Técnica} e Profissional.
Concomitantemente, também será adotado um modelo estruturado, voltado para o Jovem e seu Projeto de Vida, e assim, pretende -se elevar os índices de desempenho internos contribuindo para com a melhoria dos índices externos.
Concernente à \textbf{oferta} de Educação \textbf{Técnica} e \textbf{Profissional} uma das reformulações realizadas pelo estado de Rondônia para possibilitar aos jovens o acesso, permanência e sucesso escolar com aprendizagem dos saberes culturais e científicos para o exercício da cidadania e inserção no mundo do trabalho, foi criado o Instituto Estadual de Desenvolvimento da Educação Profissional – IDEP, por meio da Lei Complementar nº 908/2016.
Esta autarquia é vinculada à Secretaria de Estado da Educação, dotada de autonomia administrativa, pedagógica, disciplinar, financeira, orçamentária e patrimonial, e atua como \textbf{órgão} gestor da Política de Educação Profissional do Estado de Rondônia, oferecendo à sociedade rondoniense oportunidades educacionais gratuitas que visam o \textbf{fortalecimento} da capacidade \textbf{profissional} produtiva para o desenvolvimento competitivo dos estudantes no mundo do trabalho, \textbf{atendendo} as \textbf{políticas} \textbf{socioeducacionais}.
Tais medidas têm por objetivo assegurar que todos tenham acesso ao Ensino Médio de \textbf{qualidade}, propiciar o domínio das competências e habilidades das áreas de conhecimento previstas para cada etapa escolar, bem como desenvolver valores necessários para o século XXI, que contribuam para a formação de jovens competentes, dotados de autonomia e solidariedade.
Nessa perspectiva, a implementação do novo ensino médio visa incentivar o estudante a desenvolver diferentes formas de expressão, não só em âmbito intelectual, como também nas esferas artística, física, cultural, digital e social, permitindo aos estudantes percorrer \textbf{itinerários} diversificados e que melhor respondam à heterogeneidade de suas condições, interesses e aspirações, com previsão de espaços e tempos diversificados.
Mediante a esta nova abordagem metodológica, espera -se que os jovens desenvolvam o prazer pelo conhecimento e sejam estimulados a concluir os estudos nesta etapa da educação básica com êxito e perspectivas para o futuro.

% matched lemmas: atender, carga, censo, currículo, distrito, equânime, estadual, formativo, fortalecimento, horário, itinerário, média, oferta, ofertar, política, profissional, qualidade, rede, reforma, reprovação, socioeducacional, trajetória, técnico, órgão
\end{textsample}
